\documentclass[11pt]{article}
\usepackage[margin=1in]{geometry}
\usepackage{amsmath}
\usepackage{booktabs}

\title{Boosted Trees for Smooth Non-Linear Regression: The Friedman \#1 Benchmark}
\author{Paper2Proof canary set\thanks{A small paper written for the Paper2Proof evaluation set (eval set A).}}
\date{2026}

\begin{document}
\maketitle

\begin{abstract}
Friedman's first regression problem mixes a sine interaction, a quadratic term and linear terms
with five irrelevant inputs. Gradient-boosted regression trees reach a test RMSE of 1.300 and
reduce the RMSE of ordinary least squares by 1.164 on the same splits.
\end{abstract}

\section{Introduction}
Synthetic regression benchmarks with a known generating function make it easy to separate model
error from noise. We report reference errors for boosted trees on Friedman \#1.

\section{Method}
\paragraph{Boosted trees.} Gradient-boosted regression trees with squared-error loss, 100
stages, learning rate 0.1, maximum depth 3, no subsampling, random state equal to the seed.
\paragraph{Baseline.} Ordinary least squares linear regression with an intercept.

\section{Experimental setup}
\paragraph{Data.} One fixed sample of 2000 points from scikit-learn's \texttt{make\_friedman1}
with 10 input features and Gaussian noise of standard deviation 1.0, generated with
\texttt{random\_state=0}.
\paragraph{Splits.} For each seed $s \in \{0, 1, 2\}$, a random (unstratified) split with 30\,\%
test points (\texttt{train\_test\_split(test\_size=0.3, random\_state=s)}); both models use the
same split.
\paragraph{Metric.} Root mean squared error (RMSE) on the test split, lower is better; mean and
sample standard deviation over the three seeds. The reduction is the per-seed difference between
the two models, averaged over seeds.
\paragraph{Software.} Python 3.12, scikit-learn 1.9.1, CPU only.

\section{Results}
Table~\ref{tab:friedman} reports the errors. Gradient-boosted trees reach a test RMSE of
1.300 $\pm$ 0.015, and they reduce the test RMSE of linear regression by 1.164 on average over
the three seeds.

\begin{table}[t]
\centering
\caption{Test RMSE on Friedman \#1, noise 1.0 (mean $\pm$ standard deviation over seeds 0, 1, 2;
lower is better).}
\label{tab:friedman}
\begin{tabular}{lc}
\toprule
Model & RMSE \\
\midrule
Gradient-boosted trees (100 stages, depth 3) & $1.300 \pm 0.015$ \\
\midrule
Reduction with respect to linear regression & $-1.164$ \\
\bottomrule
\end{tabular}
\end{table}

\section{Conclusion}
Boosted trees roughly halve the error of a linear model on Friedman \#1.

\end{document}
